\begin{figure}[htbp]
\centering
\begin{tikzpicture}[vertex/.style={circle,fill=black,inner sep=1.7pt},>=Stealth]
\node[vertex,label=above:$p$] (p) at (0,1.2) {};
\node[vertex,label=above:$\rho(q)$] (r) at (1.4,1.2) {};
\node[vertex,label=above:$q$] (q) at (4.8,1.2) {};
\node[vertex,label=below:$p'$] (pp) at (0,-1.2) {};
\node[vertex,label=below:$\rho(q')$] (rr) at (1.4,-1.2) {};
\node[vertex,label=below:$q'$] (qq) at (4.8,-1.2) {};
\draw[dashed] (p)--(r) (pp)--(rr);
\draw[very thick,->] (r)--(q);
\draw[very thick,->] (rr)--(qq);
\draw[very thick] (r) to[bend right=55] node[left,xshift=-3pt] {$d$ new vertices} (rr);
\node at (3.3,.78) {$a$ interior vertices};
\node at (3.3,-.78) {$b$ interior vertices};
\node at (3.2,0) {new path: $q\longleftrightarrow q'$};
\end{tikzpicture}
\caption{Representative replacement when both targets have incoming
paths. Each representative is the \emph{first} interior vertex on its
incoming path. Deleting the initial edge of each incoming path leaves its
full representative-to-target segment available. The new
path has amount $a+b+d$, so the entire system has amount $L+d$.
Lemma~\ref{grow:replacement} also treats a target at a chain start and
two targets on the same original chain. Dashed initial steps are not
used by the new path.}
\label{fig:representatives}
\end{figure}
